\documentclass[11pt]{article}
\usepackage[margin=1in]{geometry}
\usepackage{amsmath,amssymb,amsthm}
\usepackage[utf8]{inputenc}
\usepackage{listings}
\usepackage{hyperref}
\usepackage{xcolor}
\usepackage{textcomp}

\definecolor{leanblue}{rgb}{0.2,0.4,0.7}

\lstdefinelanguage{Lean}{
  keywords={theorem,lemma,def,axiom,structure,class,instance,where,by,sorry,simp,exact,intro,apply,have,let,fun,match,if,then,else,import,open,namespace,end,section,variable,inductive,noncomputable,abbrev,deriving},
  keywordstyle=\color{leanblue}\bfseries,
  comment=[l]{--},
  morecomment=[s]{/-}{-/},
  string=[b]",
  basicstyle=\ttfamily\small,
  breaklines=true,
  extendedchars=true,
  inputencoding=utf8,
  literate=
    {≥}{{$\geq$}}1
    {≤}{{$\leq$}}1
    {→}{{$\to$}}1
    {←}{{$\leftarrow$}}1
    {∧}{{$\wedge$}}1
    {∨}{{$\vee$}}1
    {¬}{{$\neg$}}1
    {∀}{{$\forall$}}1
    {∃}{{$\exists$}}1
    {⊆}{{$\subseteq$}}1
    {⊂}{{$\subset$}}1
    {∈}{{$\in$}}1
    {∉}{{$\notin$}}1
    {×}{{$\times$}}1
    {⊗}{{$\otimes$}}1
    {▸}{{$\blacktriangleright$}}1
    {⟨}{{$\langle$}}1
    {⟩}{{$\rangle$}}1
    {λ}{{$\lambda$}}1
    {α}{{$\alpha$}}1
    {β}{{$\beta$}}1
    {σ}{{$\sigma$}}1
    {θ}{{$\theta$}}1
    {ε}{{$\varepsilon$}}1
    {μ}{{$\mu$}}1
    {π}{{$\pi$}}1
    {Σ}{{$\Sigma$}}1
    {Π}{{$\Pi$}}1
    {▶}{{$\triangleright$}}1
    {⊔}{{$\sqcup$}}1
    {⊓}{{$\sqcap$}}1
    {⊥}{{$\bot$}}1
    {⊤}{{$\top$}}1
    {∅}{{$\emptyset$}}1
    {∪}{{$\cup$}}1
    {∩}{{$\cap$}}1
    {≠}{{$\neq$}}1
    {ℕ}{{$\mathbb{N}$}}1
    {₀}{{$_0$}}1
    {₁}{{$_1$}}1
    {₂}{{$_2$}}1
    {|>}{{$|{\!>}$}}2
    {·}{{$\cdot$}}1
    {—}{{---}}1
    {∘}{{$\circ$}}1
    {√}{{$\sqrt{}$}}1
    {∝}{{$\propto$}}1
    {≈}{{$\approx$}}1
    {═}{{=}}1
    {⟹}{{$\Longrightarrow$}}1,
}

\newtheorem{theorem}{Theorem}
\newtheorem{lemma}[theorem]{Lemma}
\newtheorem{definition}[theorem]{Definition}
\newtheorem{axiom_stmt}[theorem]{Axiom}

\title{Formal Verification: End-to-End Finality Composition}
\author{Zach Kelling \and Woo Bin}
\date{December 2025}

\begin{document}
\maketitle

\begin{abstract}
We tie together BFT, Safety, Liveness, and Quasar into a single finality guarantee: under BFT assumptions, consensus finalizes a unique value within bounded time with hybrid (classical + post-quantum) security.
\end{abstract}

\section{Introduction}
This module composes the individual consensus results into a single end-to-end guarantee. The composition follows: BFT threshold $\to$ Safety (unique finalization) $\to$ Liveness (bounded time) $\to$ Quasar (quantum-resistant certificates). Together they establish that Lux achieves unique, bounded-time, quantum-safe finality.

\section{Definitions}
\begin{definition}[FinalityParams]
Parameters: $n$ validators, $f$ max Byzantine, $\beta$ confidence threshold, $\delta$ message delay, with $3f < n$.
\end{definition}

\begin{definition}[FinalizedBlock]
A finalized block with height, hash, BLS certificate validity, Ringtail certificate validity, and round.
\end{definition}


\section{Theorems and Axioms}
\begin{axiom_stmt}[\texttt{safety\_unique\_hash}]
Two finalized blocks at the same height must have the same hash (from Safety.safety and quorum intersection).
\end{axiom_stmt}

\begin{theorem}[\texttt{unique\_per\_height}]
Two finalized blocks at the same height have identical hashes.
\end{theorem}
\begin{proof}
Direct application of \texttt{safety\_unique\_hash}.
\end{proof}

\begin{theorem}[\texttt{bounded\_finality}]
Finality round $\leq \text{GST} + \beta + \delta$.
\end{theorem}

\begin{theorem}[\texttt{hybrid\_required}]
A finalized block requires both BLS and Ringtail certificates to be valid.
\end{theorem}
\begin{proof}
Decompose the Boolean conjunction.
\end{proof}

\begin{theorem}[\texttt{expected\_finality\_time}]
Expected finality time is $O(\beta)$ rounds: $\beta + \delta \geq \beta$.
\end{theorem}


\section{Lean Source}
\lstinputlisting[language=Lean]{../../formal/lean/Consensus/Finality.lean}

\section{References}
\noindent Source code: \url{https://github.com/luxfi/proofs/blob/main/lean/Consensus/Finality.lean}\\
\noindent White papers: \url{https://papers.lux.network}

\noindent Related white papers:
\begin{itemize}
  \item \texttt{lux-consensus.tex}
  \item \texttt{lux-quasar-consensus.tex}
\end{itemize}

\end{document}
